\editfromrepo{base}

\replacerange{The mail fraud statute}{indictment must be denied.}{[The relevant
mail fraud statute makes it a crime to use the Postal Service for various forms
of fraud and counterfeiting. The court explains that, under the case
\emph{McNally v.~United States}, 483 U.S. 350 (1987), the Supreme Court has
interpreted this mail fraud statute such that it applies only to defrauding
someone out of \emph{property}, as opposed to defrauding the public out of a
properly functioning political process.

So the defendants are liable under this rule only if they
defrauded New York City, as a property holder, out of its property. The court
first observes that the city lost license renewal fees, which are money and thus
the city's property. That's enough to overcome the motion to dismiss. But
the court goes on:]}

\replacerange{(e.g., the fee simple determinable}{and the power of
termination)}{[\having{intro-estates}{as you learned in \mref{intro-estates}}{as
you will learn in \mref{intro-estates}}{referring to the arcane topic of
``estates and future interests'' from property law}]}

\replacerange{wire fraud statutes.\readingfootnote}{Specifically,}{wire fraud
statutes. Specifically,}

\replacerange{equated with easements.}{such an easement is rare. }{equated with
easements\having{intro-easements}{.\ldots}{ [for reasons that will make more
sense after \mref{intro-easements}].}{ [for reasons beyond the scope of this
book].}\par}

\endedit

